\section{Implementation in the BX3 Framework}

The Bailout Protocol is not a standalone safety feature layered on top of an existing agent architecture. It is the runtime expression of an upstream accountability guarantee that is structurally embedded in the BX3 Framework's three-layer design. To understand how the protocol operates, one must first understand its integration points: the Purpose Layer that anchors human accountability, the Bounds Engine that detects unresolvable states and triggers escalation, and the Fact Layer that enforces the forensic ledger and hard-blocks any execution path that has not been authorized. Together, these three layers constitute the complete bail-out path from exception detection through human resolution to immutable record.

\subsection{Purpose Layer as the Human Accountability Anchor}

The BX3 Framework's Layer 1, the Purpose Layer, is the architectural locus of human accountability. Unlike architectures that assign human oversight as a configurable parameter, BX3 makes the Purpose Layer a structural requirement: it is the only layer that can originate authoritative directives, and it is the only layer that can receive an unfiltered escalation from any descendant node in the recursive tree. The Human Accountability Anchor (HAA) — the individual or governance process occupying Layer 1 — operates with two fundamental properties that make the Bailout Protocol structurally enforceable.

First, the HAA possesses \textit{irrevocable origination authority}. No execution command may enter the Fact Layer without a Purpose-layer directive authorizing it. This is not a permission check — it is a data-flow constraint. The Fact Layer will not issue an actuation signal unless the command carries a valid Purpose-layer signature. This means that when a node in the tree fires a bail-out trigger, it is not requesting permission to continue; it is reporting an inability to originate the authorization itself. The trigger does not ask: \emph{should I proceed?} It asks: \emph{I cannot authorize this myself — who can?}

Second, the HAA maintains a \textit{directional accountability bond} with every descendant node through the recursive spawning mechanism. When a parent node spawns a child by provisioning a Worksheet, the Worksheet contains a hard-coded pointer to the originating Purpose layer. This pointer is non-overrideable by any intermediate machine actor. It survives network disconnection, node failure, and protocol reassignment. A sensor in a field that has lost contact with the cloud still knows who its Human Root is — and therefore knows where to send a bail-out trigger when it encounters a condition it cannot resolve.

The Purpose Layer's primary interface to the Bailout Protocol is the Root Dashboard, which surfaces Trigger Packages from any descendant node in the tree. Each Trigger Package arriving at the HAA carries the full operational state of the originating node at the moment of trigger firing — all nine planes of the 9-Plane DAP (Mandate, Intent, Plan, Reason Log, Decision, Projection, Outcome Record, Execution, Projection Confirmation) — along with the propagation chain showing which intermediate nodes handled the exception and which escalated it. This gives the HAA complete situational awareness without requiring the human to reconstruct the failure mode from logs or telemetry.

\subsection{Bounds Engine: Trigger Detection and Escalation Initiation}

The BX3 Framework's Layer 2, the Bounds Engine, is the heuristic reasoning layer — a limbless proposer that can simulate, evaluate, and recommend but cannot directly actuate physical systems. It is the layer responsible for detecting that a node has entered an unresolvable state and for initiating the bail-out sequence. This is a critical architectural point: the Bounds Engine does not resolve the exception. It detects the exception, packages it, and propagates it upward. Resolution belongs exclusively to the HAA.

The Bounds Engine's bail-out triggering mechanism is gated by the three conditions formalized in the protocol's specification: Capability Boundary, Safety Envelope Violation (predicted), and Accountability Boundary Violation. Each condition corresponds to a distinct failure mode in which the node's local reasoning capacity is structurally exceeded.

A \textbf{Capability Boundary} trigger fires when the node encounters a sensor reading, computational result, or execution request that falls outside its provisioned operational range — not merely outside its preference, but outside the boundaries of what its internal models can interpret. The Z-Axis Indexing system plays a key role here: when a node's multi-depth sensor readings produce a Stratification Index that exceeds defined thresholds and the node has no model for the resulting condition, the Bounds Engine recognizes that its own epistemic limits have been reached. It cannot simulate forward because it does not understand the present. The bail-out fires not because the situation is dangerous, but because the node has encountered a condition that is genuinely beyond its interpretive capacity.

A \textbf{Safety Envelope Violation (predicted)} trigger fires when the Bounds Engine's simulation layer — Plane 6, the Projection — indicates that a proposed action, if executed, would violate one or more Safety Envelope parameters. Critically, this trigger fires on prediction, not on observation. The system does not wait for physical harm to occur; it fires when the simulation shows that harm is probable within the current execution path. This is the proactive core of the Bailout Protocol: it intercepts harm before it propagates rather than after it accumulates.

An \textbf{Accountability Boundary} trigger fires when the node encounters a decision point whose legal, regulatory, or ethical implications exceed the node's provisioned authority — regardless of whether the node has the operational competence to execute the action. The distinction between operational authority and accountability authority is essential. A node may be technically capable of executing an action but not legally authorized to do so. When the Bounds Engine detects this mismatch, it triggers upward rather than proceeding.

When any of the three conditions fires, the Bounds Engine assembles a Trigger Package (see Section 3, Table 1) and transmits it asynchronously to the parent node via the Recursive Spawning uplink. The transmission is fire-and-forget with delivery confirmation. The Bounds Engine does not wait for the parent to acknowledge before continuing its other operations — the trigger propagation is independent of the node's ongoing reasoning workload.

Critically, the Bounds Engine cannot suppress a trigger once it has fired. The Machine Actor Exclusion property ensures that if the propagation chain reaches an intermediate machine node — a Bounds Engine that is not the Human Root — the trigger does not stop. It continues upward. This prevents a compromised or misconfigured intermediate node from suppressing a legitimate bail-out, which would be the most dangerous failure mode in any accountable autonomous system.

\subsection{Forensic Ledger: Fact Layer as the Immutable Record-Keeper}

The BX3 Framework's Layer 3, the Fact Layer, is the physical execution firewall — the only layer that can actuate physical systems and the only layer that can write to the forensic ledger. The Fact Layer's relationship to the Bailout Protocol is dual: it is both the enforcer of the protocol's execution constraints and the archival system that records the entire trigger lifecycle for post-hoc audit.

Every Trigger Package that is transmitted, propagated, resolved, or dismissed is recorded by the Fact Layer in the forensic ledger. The ledger entry for a bail-out event captures the following at minimum: the unique trigger identifier and nanosecond timestamp, the originating node's full state snapshot (all nine planes), the specific trigger condition that fired, the confidence score assigned by the originating Bounds Engine, the complete propagation chain showing arrival and departure timestamps at each intermediate node, the HAA's directive upon resolution, and the physical outcome of the executed directive. This creates a complete, tamper-evident record of the exception's lifecycle from first detection through human resolution through physical execution.

The forensic ledger is not merely a log — it is a cryptographically sealed structure. Each ledger entry carries a SHA-256 hash of the previous entry, forming a chained sequence in which any retrospective modification to an earlier entry breaks the chain and is detectable in O(1) time. The ledger is append-only: no entry may be deleted, no timestamp may be modified, and no trigger resolution may be altered after the fact. This makes the Bailout Protocol auditable with mathematical certainty. Regulators, courts, or internal reviewers can reconstruct exactly what happened, when, at whose node, and what the human decided — without relying on the system's own self-reporting, which is precisely what the Bailout Protocol is designed to make unnecessary.

The Fact Layer's enforcement of the forensic ledger is architecturally inseparable from its enforcement of execution constraints. Before any directive from the HAA is executed by the Fact Layer, the Fact Layer independently verifies that: (a) the directive originates from a Purpose-layer signature, (b) the execution is within the Safety Envelope parameters, (c) the spatial resolution of the action is within the node's provisioned tier, and (d) the action does not conflict with any active regulatory constraint. Only when all four checks pass does the Fact Layer write the execution to the ledger and actuate the physical system. This independent verification means that even if the HAA issues a directive that would violate safety constraints — a theoretically possible human error — the Fact Layer hard-blocks the execution and logs the blocked attempt. The Bailout Protocol does not trust the Human Root without verification; it trusts the process that the Human Root governs.

The Fact Layer also enforces the Bailout Protocol's state machine transitions. The trigger lifecycle (TRIGGER\_FIRED $\rightarrow$ DELIVERING\_TO\_PARENT $\rightarrow$ PARENT\_EVALUATING $\rightarrow$ ESCALATING\_TO\_PARENT or RESOLVED\_LOCALLY $\rightarrow$ HUMAN\_ROOT\_REVIEWING $\rightarrow$ DIRECTIVE\_ISSUED $\rightarrow$ CLOSED) is managed as an extension of the Fact Layer's state machine. Each transition is recorded in the forensic ledger, and each transition's preconditions are independently enforced by the Fact Layer's deterministic execution logic. The state machine is not a soft protocol — it is a Fact Layer execution constraint. Any attempt to skip a state transition, suppress a trigger, or fabricate a resolution is architecturally blocked by the same physical-layer enforcement that prevents unauthorized actuation.

\subsection{The Bailout Protocol as Runtime Expression of the Accountability Guarantee}

The BX3 Framework's core accountability guarantee can be stated formally: \emph{for every physical action taken by any node in the recursive tree, there exists a human being who authorized that action through a Purpose-layer directive, and this authorization is verifiable in the forensic ledger before the action executes}. The Bailout Protocol is the runtime expression of this guarantee because it is the mechanism that activates when the normal authorization pathway is insufficient — when a node encounters a condition that cannot be authorized by any machine actor in the chain and must therefore propagate the authorization requirement upward until it reaches a human.

The protocol's relationship to the three-layer architecture is therefore not incidental — it is constitutive. The Purpose Layer defines the accountability structure that the protocol enforces. The Bounds Engine detects the conditions under which the protocol must activate. The Fact Layer records the protocol's activation and enforces the resolution pathway. Without all three layers operating in concert, the Bailout Protocol is unenforceable: a purpose layer without a bounds engine cannot detect unresolvable states; a bounds engine without a fact layer cannot record the detection; a fact layer without a purpose layer has no human to receive the escalation.

The directional guarantee enforced by the Bailout Protocol — the system fails upward into human consciousness, never downward into algorithmic chaos — is thus a direct consequence of this three-layer integration. Because the Fact Layer will not execute any action without a Purpose-layer signature, and because the Bounds Engine will propagate any unresolvable condition upward until it reaches a Purpose-layer node (bypassing all intermediate machine actors via the Machine Actor Exclusion property), the system's failure mode is structurally constrained: the failure always surfaces at the HAA, never in the execution layer. This directionality is not probabilistically desired — it is architecturally enforced by the interdependent constraints of all three layers acting in concert.

The forensic ledger ensures that this guarantee is not merely a present-state property but a historically verifiable one. Even if the system operates correctly for years — even if no bail-out is ever triggered — the ledger demonstrates that the guarantee was structurally maintained at every moment. Every execution traceable to a human authorization, every escalation captured, every resolution recorded. The Bailout Protocol's value is not only in the exceptions it handles but in the accountability architecture it makes permanently auditable.

\subsection{State Machine as Fact Layer Extension}

The Bailout Protocol's trigger lifecycle is implemented as a state machine whose transitions are enforced by the Fact Layer as deterministic execution constraints. The state machine is not a soft description of expected behavior — it is a Fact Layer execution protocol in which each state transition has preconditions that the Fact Layer independently verifies before allowing the transition to occur.

The trigger begins in the TRIGGER\_FIRED state when the Bounds Engine detects any of the three trigger conditions. The Fact Layer records the trigger's initial state and begins monitoring for the next transition. The trigger moves to DELIVERING\_TO\_PARENT when the Fact Layer confirms that the Trigger Package has been transmitted to the parent node via the Recursive Spawning uplink. The transition requires delivery confirmation — if delivery fails after a configurable number of retries, the trigger transitions to a DISCONNECTED\_ESCALATION state in which the node stores the trigger locally and attempts re-delivery via the Local Survivability pathway.

In the PARENT\_EVALUATING state, the parent node's Bounds Engine examines the Trigger Package and determines whether it can resolve the exception locally. If the parent can resolve — the proposed resolution is within its Safety Envelope and does not require HAA authorization — the trigger transitions to RESOLVED\_LOCALLY, and the Fact Layer records the resolution directive and propagates it back down the tree. If the parent cannot resolve, the trigger transitions to ESCALATING\_TO\_PARENT and the propagation continues.

The critical state is HUMAN\_ROOT\_REVIEWING, which is reached when the propagation chain has exhausted all machine actors without finding a resolution. In this state, the Fact Layer surfaces the Trigger Package to the HAA via the Root Dashboard and suspends execution of the originating node's pending action. The suspension is enforced by the Fact Layer: even if the HAA does not respond immediately, no autonomous execution proceeds until either the HAA issues a directive or a timeout condition is reached (at which point the system defaults to the safest available state per the Safety Envelope parameters).

When the HAA responds, the trigger transitions to DIRECTIVE\_ISSUED and the Fact Layer records the directive in the forensic ledger. The directive then propagates back down the tree through the PARENT nodes, each of which updates its state and executes the relevant portions of the directive. Finally, the trigger reaches the CLOSED state, and the Fact Layer records the complete lifecycle — from first trigger firing through propagation through human resolution through physical outcome — as a permanent, cryptographically sealed ledger entry.

The state machine's extension through the Fact Layer means that the Bailout Protocol's enforcement is as reliable as the Fact Layer's own execution guarantees. The Fact Layer is 100\% deterministic: it will not skip a state transition, will not fabricate a resolution, and will not execute an action without the required preconditions. This makes the Bailout Protocol the most reliably enforced accountability mechanism in the BX3 Framework — precisely because its enforcement is not distributed across the reasoning layers but concentrated in the single deterministic layer that controls physical execution.

\end{document}